\noindent\textbf{Elección del procedimiento.} La envolvente es exponencial: dos máximos separados por el periodo amortiguado cumplen $x_1/x_2=e^{\zeta\omega_nT_d}$. Con $T_d=2\pi/[\omega_n\sqrt{1-\zeta^2}]$, resulta $\delta=2\pi\zeta/\sqrt{1-\zeta^2}$. Elevar al cuadrado y aislar $\zeta$ da la expresión empleada.

\begin{enumerate}

\item Las amplitudes pueden quedarse ambas en mm porque solo aparece su cociente.

\[\fitmath{\delta=\ln\left(\frac{x_{1}}{x_{2}}\right)}\]

\[\fitmath{\delta=\ln\left(\frac{(10)}{(6)}\right)\approx 0{,}5108256\,}\]

\item Tomar la raíz positiva para amortiguamiento disipativo.

\[\fitmath{\zeta=\frac{\delta}{\sqrt{4\cdot {\pi}^{2}+{\delta}^{2}}}}\]

\[\fitmath{\zeta=\frac{(0{,}5108256238)}{\sqrt{4\cdot {\pi}^{2}+{(0{,}5108256238)}^{2}}}\approx 0{,}08103306\,}\]

\end{enumerate}
